\section{Other Related Work}

\textbf{Asynchronous geo-distributed training:} Several systems propose training over geo-distributed, preemptible, and heterogeneous resources by introducing asynchrony and reducing communication via quantization and sparsification. DiLoCo~\cite{douillard2024diloco} uses federated averaging to reduce communication, performing more local computation before synchronization. SWARM~\cite{ryabinin2023swarm}  proposes a decentralized, model-parallel approach to deal with poorly connected and unreliable devices. CocktailSGD~\cite{Wang23cocktailsgd} uses gradient compression to improve communication over low-bandwidth networks. These approaches influence training dynamics and are orthogonal to our work. \sailor employs synchronous training, which is preferred in large-scale training~\cite{wagenlander24tenplex}.


\textbf{Automatic VM selection}: CherryPick~\cite{Alipourfard17Cherrypick}, RAMBO~\cite{Li21Rambo}, and PARIS~\cite{Yadwadkar17Paris} apply Bayesian Optimization or performance modeling to recommend optimal VM types. SkyPilot~\cite{Zongheng23Skypilot} picks cost-efficient VMs across providers based on workload and user constraints. These methods treat workloads as black boxes and are not tailored for ML training. Srifty ~\cite{Luo22Srifty}, Cynthia~\cite{Zheng19Cynthia}, SpotDNN~\cite{Shang2023spotDNN}, DeepSpot~\cite{Lee17DeepSpotCloud} find optimal configurations of on-demand and spot instances for ML jobs. However, they target small data-parallel jobs, and are inadequate for large-scale training that involves various types of communication, which is our focus.

\textbf{Heterogeneous and Geo-distributed ML Inference:} Recent works have proposed systems for ML inference on heterogeneous and geo-distributed resources~\cite{Mei25Helix, Jian24Hexgen, Mao25Skyserve}. Although training and inference workloads differ significantly (e.g. in runtime and memory footprint estimation), synchronization and  communication patterns as well as planning decisions such as partitioning models across resources are relevant to both. Helix~\cite{Mei25Helix} and HexGen~\cite{Jian24Hexgen} consider both heterogeneous and geo-distributed resources for LLM inference. Helix uses a time-consuming MILP algorithm for model placement, which is not suitable for environments with high resource availability. HexGen uses a  more lightweight dynamic programming approach for splitting tensor and pipeline parallelism across resources, making it more appropriate for dynamic settings.
However, both of these works optimize only for performance, ignoring the monetary communication costs that arise in geo-distributed settings. SkyServe~\cite{Mao25Skyserve} places replicas of models across zones and regions, but restricts each replica to a single zone and homogeneous resources. As replicas do not communicate in inference, inter-zone communication is not considered. In contrast, \sailor targets geo-distributed training, which requires frequent communication among data-parallel replicas, substantially increasing scheduling complexity.
